\documentclass[11pt]{article}
\usepackage[margin=1in]{geometry}
\usepackage{amsmath,amssymb,booktabs}
\usepackage[T1]{fontenc}
\usepackage{xurl}
\usepackage[hidelinks]{hyperref}
\begin{document}
\title{Attempt 209 calibration and shape comparison: draft for transcription}
\author{}\date{}\maketitle
\noindent\textit{Draft for the researcher to adapt. Run 1 was used for calibration; Runs 2 and 3 are exploratory frozen-parameter comparisons. The historical search objective and the later corrected-convention evaluation metric are different. Editorial notes at the end are not manuscript text.}

\section{Model calibration}
\subsection{Growth parameters and repeated-hit response}
The model represents the local effect of the wheel through effective layer-growth inputs rather than independently measured constitutive parameters. The four calibrated, dimensionless inputs are $\boldsymbol\theta=(g_t,g_b,o_t,o_b)$, where $g_\ell$ and $o_\ell$ are the growth magnitude and orthotropy coefficient for layer $\ell\in\{t,b\}$. The directional inputs along and transverse to a local track are
\begin{equation}
 g_{\ell,\parallel}=g_\ell(1+o_\ell),\qquad
 g_{\ell,\perp}=g_\ell(1-o_\ell).
 \label{eq:directional-growth}
\end{equation}
A face-shared history modulates each growth increment. If $n$ previous plastic hits have been counted on a mesh face, the multiplier applied to its next hit is
\begin{equation}
 q(n)=f+(1-f)e^{-\beta n}.
 \label{eq:hit-decay}
\end{equation}
Here $\beta$ and $f$ were fixed at $0.2231435513$ and $0$, respectively; they were not fitted in this search. Thus the first hit has $q(0)=1$. This phenomenological hit-count decay is not a conventional stress--plastic-strain hardening law. The prescribed profile and track overlap also affect the local increments.

\begin{table}[ht]
\centering
\caption{Attempt 209 parameters and the final expanded search bounds. All values are dimensionless. Earlier stages used different bounds.}
\label{tab:attempt209-parameters}
\begin{tabular}{lrrr}\toprule
Parameter & Lower bound & Upper bound & Selected value \\\midrule
$g_t$ & $5\times10^{-5}$ & $10^{-2}$ & $2.88351\times10^{-4}$ \\
$g_b$ & $10^{-5}$ & $10^{-3}$ & $1.45880\times10^{-5}$ \\
$o_t$ & $-1$ & $1$ & $-0.398916$ \\
$o_b$ & $-1$ & $1$ & $0.960306$ \\\bottomrule
\end{tabular}
\end{table}

The sheet model was $200\times300\times0.77$ mm, with an effective track width of $2$ mm. Its structured $2$ mm mesh comprised 15,251 vertices and 30,000 triangles. The prescribed history was accumulated before a final unconstrained equilibrium solve. The numerical material settings used $E=1$ and $\nu=0.33$; $E=1$ is a numerical setting, not a measured modulus. The hybrid equilibrium solver used a quasi-Newton warm-up followed by sparse Newton iterations, with a gradient-norm target of $5\times10^{-14}$ in internal solver units. The mesh size alone is not a mesh-convergence study.

\subsection{Historical Run 1 optimization}
The reported optimization fitted Run 1 only. For this historical material-reference comparison, the predicted shell was framed with respect to corresponding undeformed mesh vertices and the scan was compared without optimizing a scan-to-model rigid pose. Let $\mathbf p_i$ be the retained scan points, $S(\boldsymbol\theta)$ the predicted triangular surface, $\mathbf v_j(\boldsymbol\theta)$ the predicted vertices, and $w_j$ their normalized reference dual areas. The historical objective was
\begin{align}
 d_i(\boldsymbol\theta)&=\min_{\mathbf s\in S(\boldsymbol\theta)}\lVert\mathbf p_i-\mathbf s\rVert_2,\\
 r_j(\boldsymbol\theta)&=\min_i\lVert\mathbf v_j(\boldsymbol\theta)-\mathbf p_i\rVert_2,\\
 J_{\mathrm{hist}}(\boldsymbol\theta)&=\left\{\tfrac12\left[\frac{1}{N}\sum_{i=1}^{N}d_i^2+\sum_{j=1}^{M}w_jr_j^2\right]\right\}^{1/2}.
 \label{eq:historical-objective}
\end{align}
Thus the forward term measured scan-point-to-model-surface distance, while the reverse term used distances from model vertices to the scan point cloud. Both terms are in millimetres. The reverse sampling and the scan footprint influence this historical score; it is not the later shared-footprint surface metric below.

The search used a Gaussian-process surrogate with a Mat\'ern-$5/2$ covariance for the negative logarithm of the positive objective; growth magnitudes used logarithmic coordinates and orthotropy used linear coordinates. The final continuation proposed batches of four using a success-aware expected-improvement procedure with a separate model for solver completion. This describes the verified final continuation, not necessarily every earlier proposal policy.

The reported fixed-scan sequence used 80 rescored historical records as prior observations and added 160 evaluations in successive allocations of 40, 40 and 80. Of the 160 additional records, 151 had valid objectives and nine timed out. These staged allocations ended at their planned evaluation counts, not at a parameter-convergence criterion. Attempt 209 had the lowest valid recorded historical objective, $J_{\mathrm{hist}}=1.182$ mm. These archived counts do not establish every earlier solver invocation or a global optimum.

\subsection{Corrected-convention shape evaluation}
The selected parameters were subsequently held fixed for corrected-convention forward evaluations of Runs 1--3. This later evaluation used a revised shared-footprint, double-sided surface discrepancy rather than Eq.~\eqref{eq:historical-objective}. Public coordinates are $x$ for width, $z$ for length and $y$ for height. The measured top surface is approximated as the shell midsurface; the scan and model are placed in the prescribed material-reference frame without optimizing a rigid scan-to-model pose.

For run $c$, let $\widehat{\mathbf p}_{ci}$ be its retained scan points, $\widehat S_c$ the framed simulated triangle surface, and $\widehat T_c$ a triangulated scan surface on the shared footprint. Let $\widehat{\mathbf q}_{ci}$ be the model surface sampled on that footprint and $a_{ci}$ normalized scan-triangle dual-area weights. The revised evaluation uses
\begin{align}
 d_{ci}&=\min_{\mathbf s\in\widehat S_c}\lVert\widehat{\mathbf p}_{ci}-\mathbf s\rVert_2,\\
 r_{ci}&=\min_{\mathbf t\in\widehat T_c}\lVert\widehat{\mathbf q}_{ci}-\mathbf t\rVert_2,\\
 D_c&=\left\{\tfrac12\left[\frac{1}{N_c}\sum_{i=1}^{N_c}d_{ci}^{2}+\sum_{i=1}^{N_c}a_{ci}r_{ci}^{2}\right]\right\}^{1/2}.
 \label{eq:revised-evaluation}
\end{align}
The scan triangulation excludes long edges that could bridge gaps. The prediction is sampled using its actual triangle connectivity; material-coordinate interpolation is used where vertical interpolation on the deformed projection is unavailable. Identical surfaces give zero discrepancy within numerical tolerance. This identity check does not imply that a real scan can be fitted exactly. The revised values in Table~\ref{tab:attempt209-fit} must not be plotted as though they were the historical optimization trajectory.

\begin{table}[ht]
\centering
\caption{Corrected-convention discrepancies at frozen Attempt 209 parameters, in mm. $P_{95}$ is the unweighted 95th percentile of scan-to-model distances; reverse RMS uses normalized scan-footprint area weights.}
\label{tab:attempt209-fit}
\begin{tabular}{llrrrr}\toprule
Case & Role & $D_c$ & Forward RMS & Reverse RMS & Forward $P_{95}$ \\\midrule
Run 1 & calibration case & 0.626 & 0.630 & 0.622 & 1.184 \\
Run 2 & exploratory transfer & 0.678 & 0.687 & 0.670 & 1.324 \\
Run 3 & exploratory transfer & 0.719 & 0.726 & 0.713 & 1.389 \\\bottomrule
\end{tabular}
\end{table}

The three values in Table~\ref{tab:attempt209-fit} are descriptive geometric discrepancies for one scan per case, not uncertainty estimates from repeated experiments. Runs 2 and 3 were previously inspected during development and are not pristine blind validation cases. The later Run 1 forward evaluation also used corrected solver and toolpath conventions, so its revised discrepancy should not be described as the score minimized by the historical search.

\subsection{Shape and parameter diagnostics}
The full supported height-difference maps and their empirical distributions show where the saved model and scan differ. Fixed sections at $z=-35$, $65$ and $165$ mm provide complementary width--height comparisons. The signed difference $\Delta y(x,z)=y_{\mathrm{model}}(x,z)-y_{\mathrm{experiment}}(x,z)$ is a vertical height difference, not the Euclidean surface distance in Eq.~\eqref{eq:revised-evaluation}. The three-dimensional overlays use an eightfold visual height exaggeration, and the displayed scan points are a fixed-seed sample; numerical tick values are not exaggerated.

A retrospective Gaussian process fitted to 221 valid archived observations provides conditional one-parameter slices of the historical objective. Each slice fixes the other three parameters at Attempt 209; its pointwise bands are conditional on the fitted surrogate and are not confidence intervals for physical parameters. The observations were selected adaptively. The selected $o_b$ lies near its upper search bound and $g_b$ lies toward the lower end of its logarithmic interval. These observations and the conditional slices do not establish a unique parameter set, physical identifiability or the necessity of all four parameters.

\section*{Editorial notes and evidence (not manuscript text)}
\begin{itemize}
\item Public Run 1, 2 and 3 correspond to the historical identifiers \texttt{run6}, \texttt{run9} and \texttt{run10}, respectively. Historical file names are retained for provenance.
\item The selected figures, captions and input identities are in \path{calibration/paper_figures/attempt209_publication_clean/}. They include historical optimization history and GP slices, corrected toolpaths, height residual distributions, fixed sections and three-case overlays. Follow individual captions for sampling, interpolation and display details.
\item Historical optimization evidence: \path{run/run6_expanded_bo_continue20_recovery/evaluations.jsonl}, \path{bayes_summary.json} and \path{bayes_state.json} in that directory; reconstruction and limits: \path{docs/calibration/bayesian_calibration/calibration_draft_verification.md}. Corrected case metrics: \path{run/convention_corrected/} and the figure provenance. Prior draft \path{calibration_manuscript_draft.tex} describes a different joint fit and is preserved separately.
\item Before submission, verify the physical fixture, worked face, scan-to-midsurface approximation, experimental dimensions, figure dimensions and bibliography against laboratory and journal records. Add appropriate, independently checked citations for Bayesian optimization and the numerical solver. Earlier-stage acquisition details and independent experimental replication remain limited.
\end{itemize}
\end{document}
